% GENERATED FILE. Edit canonical lessons, then run npm run generate:books.
% canonical-source-hash: fnv1a64:087b428688bcb3c1
% canonical-lessons: SA-C91-prabhatam, SA-C91-sandhya, SA-C91-hemantah, SA-C91-sisirah, SA-C91-sarad

\chapter{Morning, Twilight, Early winter, Late winter, Autumn}
\label{ch:sa-morning-twilight-early-winter-late-winter-autumn}
\hlchaptermodality{91}

\begin{chapteropening}
\textbf{By the end of this chapter:} \emph{I can say morning, twilight, early winter, late winter and autumn.}

Complete the last lesson of chapter 91: i can say morning, twilight, early winter, late winter and autumn.
\end{chapteropening}

\section[prabhātam]{\sk{प्रभातम्} (prabhātam) — morning, daybreak}

\label{lesson:SA-C91-prabhatam}



\begin{quote}
Before the new one: say the Sanskrit for midnight, then the Sanskrit for dawn.
\end{quote}

\subsection*{You'll want to know: \sk{प्रभातम्}}
\textbf{\sk{प्रभातम्}} — \emph{prabhātam} — \textquotedblleft{}morning, daybreak\textquotedblright{}.

\begin{grammarlens}[title={What you've built}]
One more word for this chapter.
\end{grammarlens}

\subsection*{Your turn}
\begin{itemize}
\raggedright
  \item \emph{Say it:} \emph{prabhātam}
  \item \emph{Say it:} \emph{prabhātam}, once more
  \item \emph{Say it:} say \emph{prabhātam}
  \item \emph{Recall:} say the Sanskrit for midnight, then the Sanskrit for dawn, then say \emph{prabhātam} again
\end{itemize}

\begin{tcolorbox}[breakable,colback=teal!4,colframe=teal!35!black,title={Before you move on}]
What does \sk{प्रभातम्} mean? (\textquotedblleft{}Morning, daybreak\textquotedblright{}.) Say it once more.
\end{tcolorbox}
\section[sandhyā]{\sk{सन्ध्या} (sandhyā) — twilight}

\label{lesson:SA-C91-sandhya}



\begin{quote}
Before the new one: say the Sanskrit for dawn, then the Sanskrit for morning.
\end{quote}

\subsection*{You'll want to know: \sk{सन्ध्या}}
\textbf{\sk{सन्ध्या}} — \emph{sandhyā} — \textquotedblleft{}twilight\textquotedblright{}.

\begin{grammarlens}[title={What you've built}]
One more word for this chapter.
\end{grammarlens}

\subsection*{Your turn}
\begin{itemize}
\raggedright
  \item \emph{Say it:} \emph{sandhyā}
  \item \emph{Say it:} \emph{sandhyā}, once more
  \item \emph{Say it:} say \emph{sandhyā}
  \item \emph{Recall:} say the Sanskrit for dawn, then the Sanskrit for morning, then say \emph{sandhyā} again
\end{itemize}

\begin{tcolorbox}[breakable,colback=teal!4,colframe=teal!35!black,title={Before you move on}]
What does \sk{सन्ध्या} mean? (\textquotedblleft{}Twilight\textquotedblright{}.) Say it once more.
\end{tcolorbox}
\section[hemantaḥ]{\sk{हेमन्तः} (hemantaḥ) — early winter}

\label{lesson:SA-C91-hemantah}



\begin{quote}
Before the new one: say the Sanskrit for morning, then the Sanskrit for twilight.
\end{quote}

\subsection*{You'll want to know: \sk{हेमन्तः}}
\textbf{\sk{हेमन्तः}} — \emph{hemantaḥ} — \textquotedblleft{}early winter\textquotedblright{}.

\begin{grammarlens}[title={What you've built}]
One more word for this chapter.
\end{grammarlens}

\subsection*{Your turn}
\begin{itemize}
\raggedright
  \item \emph{Say it:} \emph{hemantaḥ}
  \item \emph{Say it:} \emph{hemantaḥ}, once more
  \item \emph{Say it:} say \emph{hemantaḥ}
  \item \emph{Recall:} say the Sanskrit for morning, then the Sanskrit for twilight, then say \emph{hemantaḥ} again
\end{itemize}

\begin{tcolorbox}[breakable,colback=teal!4,colframe=teal!35!black,title={Before you move on}]
What does \sk{हेमन्तः} mean? (\textquotedblleft{}Early winter\textquotedblright{}.) Say it once more.
\end{tcolorbox}
\section[śiśiraḥ]{\sk{शिशिरः} (śiśiraḥ) — late winter}

\label{lesson:SA-C91-sisirah}



\begin{quote}
Before the new one: say the Sanskrit for twilight, then the Sanskrit for early winter.
\end{quote}

\subsection*{You'll want to know: \sk{शिशिरः}}
\textbf{\sk{शिशिरः}} — \emph{śiśiraḥ} — \textquotedblleft{}late winter\textquotedblright{}.

\begin{grammarlens}[title={What you've built}]
One more word for this chapter.
\end{grammarlens}

\subsection*{Your turn}
\begin{itemize}
\raggedright
  \item \emph{Say it:} \emph{śiśiraḥ}
  \item \emph{Say it:} \emph{śiśiraḥ}, once more
  \item \emph{Say it:} say \emph{śiśiraḥ}
  \item \emph{Recall:} say the Sanskrit for twilight, then the Sanskrit for early winter, then say \emph{śiśiraḥ} again
\end{itemize}

\begin{tcolorbox}[breakable,colback=teal!4,colframe=teal!35!black,title={Before you move on}]
What does \sk{शिशिरः} mean? (\textquotedblleft{}Late winter\textquotedblright{}.) Say it once more.
\end{tcolorbox}
\section[śarad]{\sk{शरद्} (śarad) — autumn}

\label{lesson:SA-C91-sarad}



\begin{quote}
Before the new one: say the Sanskrit for early winter, then the Sanskrit for late winter.
\end{quote}

\subsection*{You'll want to know: \sk{शरद्}}
\textbf{\sk{शरद्}} — \emph{śarad} — \textquotedblleft{}autumn\textquotedblright{}.

\begin{grammarlens}[title={What you've built}]
One more word for this chapter.
\end{grammarlens}

\subsection*{Your turn}
\begin{itemize}
\raggedright
  \item \emph{Say it:} \emph{śarad}
  \item \emph{Say it:} \emph{śarad}, once more
  \item \emph{Say it:} say \emph{śarad}
  \item \emph{Recall:} say the Sanskrit for early winter, then the Sanskrit for late winter, then say \emph{śarad} again
\end{itemize}

\begin{tcolorbox}[breakable,colback=teal!4,colframe=teal!35!black,title={Before you move on}]
What does \sk{शरद्} mean? (\textquotedblleft{}Autumn\textquotedblright{}.) Say it once more.
\end{tcolorbox}
